\documentclass{article}
\usepackage[utf8]{inputenc}
\usepackage[T1]{fontenc}
\usepackage[french]{babel}
\usepackage{amsmath}   % Pour les environnements d'équations et \text
\usepackage{mathtools} % Pour mathtools (equation*)
\usepackage{amssymb}   % Pour les symboles avancés (\nleq, \ngeq, etc.)

\begin{document}

\section{Introduction aux mathématiques en LaTeX}

\subsection{Mode normal et mode math}
Écriture en mode normal : 3-2+4=5. \\
Écriture en mode math : $3-2+4=5$.

Comparaison avec l'exposant (accent circonflecte \verb|^|) :
En mode normal : 2^{10}=1024 (l'accent est affiché seul ou décale le texte).
En mode math : $2^{10} = 1024$.

\subsection{Indices et produits}
Variables indicées avec le symbole \verb|_| et \ldots : 
$a_1, a_2, \dots, a_n$.

Produit des variables avec \verb|\times| et \verb|\cdots| :
$a_1 \times a_2 \times \cdots \times a_n$.

\subsection{Inégalités, relations et fonctions}
Inégalités strictes et larges : 
$3 < 4$, $5 > 4$, $3 \leq 4$ (\verb|\leq|) et $5 \geq 4$ (\verb|\geq|).

Quantificateurs : $\exists$ (\verb|\exists|) et $\forall$ (\verb|\forall|).

Relations diverses : 
$a \neq b$ (\verb|\neq|), 
$a \equiv b$ (\verb|\equiv|), 
$P \iff Q$ (\verb|\iff|), 
$x \sim y$ (\verb|\sim|), 
$f \hookrightarrow G$ (\verb|\hookrightarrow|).

Fonctions et opérateurs : 
$\cos(x)$, $\sin(x)$, $\tan(x)$, $\lim_{x \to 0} f(x)$.

Fraction et constante : 
$\frac{a}{b}$ et l'infini $\infty$ (\verb|\infty|).

Symboles avancés (paquet \texttt{amssymb}) :
$3 \nleq 4$ (\verb|\nleq|), $3 \ngeq 4$ (\verb|\ngeq|), et $A \implies B$ (\verb|\implies|).

\section{Équations sur une ligne à part}

Équation non numérotée avec \texttt{equation*} :
\begin{equation*}
    e = mc^2
\end{equation*}

Équation numérotée avec \texttt{equation} et étiquette (\verb|\label|) :
\begin{equation}
    \label{eq:einstein}
    E = mc^2
\end{equation}

Référence à l'équation précédente avec \verb|\eqref{eq:einstein}| : 
Comme on le voit dans l'équation \eqref{eq:einstein}, l'énergie est liée à la masse. 
(Utiliser \verb|\ref| donne simplement le numéro, tandis que \verb|\eqref| ajoute les parenthèses automatiques).

\subsection{Texte et formules complexes}

Formule avec du texte intégré via \verb|\text| :
\begin{equation*}
    \forall k < n, \text{ $k$ est un diviseur de } n!
\end{equation*}

Limite complexe avec les commandes \verb|\to|, \verb|\theta|, et \verb|\pi| :
\begin{equation*}
    \lim_{\theta \to \left(\frac{\pi}{2}\right)^-} \tan(\theta) = \infty
\end{equation*}

\end{document}